% ================================================================
%  sesion19.tex  —  Sesión 19 de 20
%  Ensayo y mejora
%  Bloque 4: Producto final · Matemáticas 1.º ESO
% ================================================================
\documentclass[aspectratio=169, 12pt, spanish]{beamer}
\usepackage{almeria-beamer}

\renewcommand{\SessionNumber}{19}
\renewcommand{\SessionTitle}{Ensayo y mejora}
\renewcommand{\SessionName}{Sesión 19}
\renewcommand{\SessionObjective}{%
  Revisar y mejorar la guía turística matemática mediante
  autoevaluación y coevaluación estructurada, identificando
  errores conceptuales y de presentación.}
\renewcommand{\BlockName}{Bloque 4: Producto final}

\begin{document}

\begin{frame}[plain]\titlepage\end{frame}

\begin{frame}{Índice}
  \begin{columns}[t]
    \begin{column}{0.5\textwidth}
      \begin{enumerate}
        \item ¿Por qué revisar?
        \item Tipos de errores a buscar
        \item Rúbrica de autoevaluación
        \item Protocolo de coevaluación
        \item Errores más comunes del proyecto
      \end{enumerate}
    \end{column}
    \begin{column}{0.5\textwidth}
      \begin{enumerate}\setcounter{enumi}{5}
        \item El ciclo de mejora
        \item Ejercicios: evaluar y corregir
        \item Evidencia del día
        \item Error de actitud
        \item Resumen
      \end{enumerate}
    \end{column}
  \end{columns}
\end{frame}

% ---------------------------------------------------------------
\seccionframe{1. ¿Por qué revisar?}
% ---------------------------------------------------------------

\begin{frame}{La revisión es parte del proceso matemático}
  \begin{columns}[c]
    \begin{column}{0.52\textwidth}
      \begin{teoriabox}
        Ningún trabajo matemático se entrega sin revisión.
        Los matemáticos, ingenieros y científicos revisan
        sus trabajos \textbf{varias veces} antes de publicarlos.\\[6pt]
        El primer borrador \textbf{siempre} contiene:
        \begin{itemize}
          \item Errores de cálculo
          \item Unidades incorrectas
          \item Gráficas incompletas
          \item Datos sin verificar
          \item Explicaciones poco claras
        \end{itemize}
      \end{teoriabox}
    \end{column}
    \begin{column}{0.44\textwidth}
      \begin{tikzpicture}[
        nodo/.style={
          rounded corners=5pt,
          draw=AlmAzulM, fill=AlmFondo,
          font=\small\bfseries,
          text width=2.2cm, align=center,
          inner sep=5pt
        },
        flecha/.style={->, AlmAmbar, line width=1.5pt, shorten >=2pt}
      ]
        \node[nodo, fill=BCrear!15] (c) at (0,3.5) {Crear};
        \node[nodo] (r) at (2.5,2)   {Revisar};
        \node[nodo] (i) at (2.5,0)   {Identificar\\error};
        \node[nodo] (co)at (0,-1.2)  {Corregir};
        \node[nodo] (v) at (-2.5,0)  {Verificar};
        \node[nodo, fill=AlmVerde!15] (p) at (-2.5,2) {Presentar};
        \draw[flecha] (c)  -- (r);
        \draw[flecha] (r)  -- (i);
        \draw[flecha] (i)  -- (co);
        \draw[flecha] (co) -- (v);
        \draw[flecha] (v)  -- (p);
        \draw[flecha, dashed] (r) to[bend left=20] (c)
          node[midway, above right, font=\tiny, color=AlmGris]
          {¿Todo bien?};
      \end{tikzpicture}
    \end{column}
  \end{columns}
\end{frame}

% ---------------------------------------------------------------
\seccionframe{2. Tipos de errores}
% ---------------------------------------------------------------

\begin{frame}{Cuatro tipos de errores a buscar}
  \begin{columns}[T]
    \begin{column}{0.48\textwidth}
      \begin{tcolorbox}[enhanced, colback=BRecordar!8,
        colframe=BRecordar, fonttitle=\bfseries\small\color{AlmBlanco},
        title={\faCalculator\enskip Errores matemáticos},
        attach boxed title to top left={yshift=-2.2mm, xshift=4mm},
        boxed title style={colback=BRecordar, arc=3pt},
        arc=4pt, leftrule=4pt]
        \footnotesize
        \begin{itemize}
          \item Fórmula incorrecta ($A = b + h$ en vez de $b \times h$)
          \item Operación aritmética errónea
          \item Unidad equivocada (m² en un perímetro)
          \item Raíz cuadrada olvidada en Pitágoras
          \item Coordenadas invertidas $(y,x)$ en vez de $(x,y)$
        \end{itemize}
      \end{tcolorbox}
      \vspace{6pt}
      \begin{tcolorbox}[enhanced, colback=BAnalizar!8,
        colframe=BAnalizar, fonttitle=\bfseries\small\color{AlmBlanco},
        title={\faDatabase\enskip Errores de datos},
        attach boxed title to top left={yshift=-2.2mm, xshift=4mm},
        boxed title style={colback=BAnalizar, arc=3pt},
        arc=4pt, leftrule=4pt]
        \footnotesize
        \begin{itemize}
          \item Dato inventado sin fuente
          \item Dato obsoleto (más de 5 años)
          \item Dato sin unidades
          \item Valor implausible (Alcazaba 430\,m²)
        \end{itemize}
      \end{tcolorbox}
    \end{column}
    \begin{column}{0.48\textwidth}
      \begin{tcolorbox}[enhanced, colback=BComprender!8,
        colframe=BComprender, fonttitle=\bfseries\small\color{AlmBlanco},
        title={\faChartBar\enskip Errores de representación},
        attach boxed title to top left={yshift=-2.2mm, xshift=4mm},
        boxed title style={colback=BComprender, arc=3pt},
        arc=4pt, leftrule=4pt]
        \footnotesize
        \begin{itemize}
          \item Ejes sin etiqueta o sin unidades
          \item Escala no regular (saltos irregulares)
          \item Eje $y$ no empieza en 0 sin justificación
          \item Tipo de gráfica inadecuado
          \item Tabla no ordenada por $x$
        \end{itemize}
      \end{tcolorbox}
      \vspace{6pt}
      \begin{tcolorbox}[enhanced, colback=BEvaluar!8,
        colframe=BEvaluar, fonttitle=\bfseries\small\color{AlmBlanco},
        title={\faCommentAlt\enskip Errores de comunicación},
        attach boxed title to top left={yshift=-2.2mm, xshift=4mm},
        boxed title style={colback=BEvaluar, arc=3pt},
        arc=4pt, leftrule=4pt]
        \footnotesize
        \begin{itemize}
          \item Descripción sin datos numéricos
          \item Cálculo sin explicación de los pasos
          \item Conclusión vaga («hay muchos turistas»)
          \item Vocabulario matemático incorrecto
        \end{itemize}
      \end{tcolorbox}
    \end{column}
  \end{columns}
\end{frame}

% ---------------------------------------------------------------
\seccionframe{3. Rúbrica de autoevaluación}
% ---------------------------------------------------------------

\begin{frame}{Rúbrica de autoevaluación}
  \begin{center}
  \footnotesize
  \begin{tabular}{p{2.8cm} p{2.6cm} p{2.6cm} p{2.6cm} p{2.6cm}}
    \toprule
    \textbf{Criterio}
      & \textbf{4 — Excelente}
      & \textbf{3 — Bien}
      & \textbf{2 — Regular}
      & \textbf{1 — Mejora}\\
    \midrule
    Exactitud\\matemática
      & Todos los cálculos correctos
      & 1--2 errores menores
      & 3--4 errores
      & Errores graves\\[4pt]
    Unidades
      & Siempre correctas
      & Casi siempre
      & Algunas mal
      & Frecuentemente mal\\[4pt]
    Gráficas
      & Completas y claras
      & Falta 1 elemento
      & Faltan 2 elementos
      & Incompletas\\[4pt]
    Selección\\de datos
      & Verificados y relevantes
      & Mayoría correctos
      & Algunos irrelevantes
      & Sin criterio\\[4pt]
    Comunicación
      & Clara y matemática
      & Casi siempre clara
      & A veces confusa
      & Vaga o incorrecta\\
    \bottomrule
  \end{tabular}
  \end{center}
  \vspace{4pt}
  \begin{center}
    \small\textbf{Puntuación máxima: 20 puntos}\\
    {\footnotesize 18--20: Sobresaliente · 14--17: Notable ·
    10--13: Bien · $<$10: Necesita revisión}
  \end{center}
\end{frame}

% ---------------------------------------------------------------
\seccionframe{4. Protocolo de coevaluación}
% ---------------------------------------------------------------

\begin{frame}{Coevaluación: revisión entre equipos}
  \begin{columns}[c]
    \begin{column}{0.55\textwidth}
      \begin{teoriabox}
        \textbf{Protocolo «Estrella + Deseo»}\\[6pt]
        \begin{enumerate}
          \item \textbf{Intercambia} la guía con otro equipo.
          \item Lee con atención \textbf{3 secciones concretas}.
          \item Para cada sección anota:\\[4pt]
            \fastar\ \textbf{Estrella} (1): algo que funciona muy bien\\
            \fastar\ \textbf{Estrella} (2): otro punto positivo\\
            \faCommentAlt\ \textbf{Deseo} (1): algo a mejorar\\
            (sé específico: «la tabla de la pág.~3
            no tiene unidades en la cabecera»)
          \item Devuelve la guía con tus notas escritas.
          \item Lee el feedback recibido y \textbf{decide}
                qué cambios aplicar.
        \end{enumerate}
      \end{teoriabox}
    \end{column}
    \begin{column}{0.41\textwidth}
      \textbf{Normas de la coevaluación:}
      \begin{itemize}
        \item Sé \textbf{específico}: no «está bien» sino
              «la gráfica de barras tiene los ejes
              correctamente etiquetados con unidades».
        \item Sé \textbf{constructivo}: «En la pág.~4 falta
              la unidad del área. Añade m² al resultado
              para que sea correcto.»
        \item No es una crítica personal: estás
              evaluando el \textbf{trabajo matemático},
              no a la persona.
        \item El receptor del feedback puede
              \textbf{aceptar o rechazar} cada sugerencia
              con justificación matemática.
      \end{itemize}
    \end{column}
  \end{columns}
\end{frame}

% ---------------------------------------------------------------
\seccionframe{5. Errores más comunes}
% ---------------------------------------------------------------

\begin{frame}{Los 8 errores más frecuentes en este proyecto}
  \begin{columns}[T]
    \begin{column}{0.48\textwidth}
      \begin{enumerate}
        \item \incorrecto\ Perímetro en m²\\
              \correcto\ Perímetro en m (lineal)
        \vspace{4pt}
        \item \incorrecto\ Coordenadas como $(y,x)$\\
              \correcto\ Siempre $(x,y)$: primero horizontal
        \vspace{4pt}
        \item \incorrecto\ Diagrama de barras para temperatura\\
              \correcto\ Gráfica lineal para variables continuas
        \vspace{4pt}
        \item \incorrecto\ Tabla desordenada\\
              \correcto\ Ordenada por la variable $x$
      \end{enumerate}
    \end{column}
    \begin{column}{0.48\textwidth}
      \begin{enumerate}\setcounter{enumi}{4}
        \item \incorrecto\ $d = 3^2 + 4^2 = 25$ (falta $\sqrt{}$)\\
              \correcto\ $d = \sqrt{3^2+4^2} = 5$
        \vspace{4pt}
        \item \incorrecto\ «La Alcazaba tiene 43.000»\\
              \correcto\ «\SI{43000}{m^2}» (siempre con unidad)
        \vspace{4pt}
        \item \incorrecto\ Eje $y$ empieza en 2.000.000\\
              \correcto\ Empieza en 0 (o justifica por qué no)
        \vspace{4pt}
        \item \incorrecto\ $A = b + h$\\
              \correcto\ $A = b \times h$ para el rectángulo
      \end{enumerate}
    \end{column}
  \end{columns}
\end{frame}

% ---------------------------------------------------------------
\seccionframe{7. Ejercicios: evaluar y corregir}
% ---------------------------------------------------------------

\begin{frame}{Ejercicios multinivel}
  \begin{columns}[T]
    \begin{column}{0.48\textwidth}
      \begin{evaluarbox}
        \footnotesize\dificultad{3}\\[4pt]
        Aplica la rúbrica de autoevaluación
        a \textbf{tu propia guía}.\\[4pt]
        \begin{enumerate}
          \item Puntúa cada criterio (1--4).
          \item Calcula la puntuación total.
          \item Identifica tus \textbf{3 puntos más débiles}.
          \item Escribe un plan concreto de mejora
                para cada punto: qué cambiarás y cómo.
        \end{enumerate}
      \end{evaluarbox}
    \end{column}
    \begin{column}{0.48\textwidth}
      \begin{crearbox}
        \footnotesize\dificultad{4}\\[4pt]
        Basándote en el \textbf{feedback del equipo revisor},
        aplica \textbf{3 mejoras concretas} a tu guía.\\[4pt]
        Para cada mejora documenta:
        \begin{enumerate}
          \item \textbf{Versión original} (copia textual)
          \item \textbf{Problema identificado}
                (tipo de error y criterio de la rúbrica)
          \item \textbf{Versión corregida}
                (con el cambio aplicado)
        \end{enumerate}
        Adjunta la documentación de mejoras
        a tu guía final.
      \end{crearbox}
    \end{column}
  \end{columns}
\end{frame}

% ---------------------------------------------------------------
\begin{frame}{Evidencia del día}
  \begin{evidenciabox}
    Entrega:
    \begin{enumerate}
      \item \textbf{Rúbrica de autoevaluación} completada
            con puntuación numérica en cada criterio y
            reflexión escrita sobre los 3 puntos más débiles.
      \item \textbf{Guía revisada} con las correcciones
            aplicadas marcadas claramente (usa un color diferente
            para las correcciones).
      \item \textbf{Hoja de mejoras} documentando las 3 correcciones
            (versión original → problema → versión corregida).
    \end{enumerate}
    \vspace{4pt}
    \textit{Una guía bien revisada vale mucho más que una
    guía apresurada con muchos errores. La calidad matemática
    es lo que hace profesional vuestro trabajo.}
  \end{evidenciabox}
\end{frame}

% ---------------------------------------------------------------
\begin{frame}{¡Cuidado con este error de actitud!}
  \begin{errorbox}
    \textbf{«Mi cálculo está bien porque yo lo hice»}\\[6pt]
    \incorrecto\ Resistirse al feedback y defender los errores
    por orgullo. La revisión crítica es debilidad.\\[8pt]
    \correcto\ Los mejores matemáticos y científicos del mundo
    \textbf{dan la bienvenida a las correcciones} porque les
    permiten mejorar. En matemáticas hay una respuesta correcta
    y es \textbf{verificable}. Si alguien señala un error, es
    una oportunidad, no un ataque.\\[6pt]
    «\textit{Un error descubierto hoy es mejor que
    un error entregado mañana.}»
  \end{errorbox}
\end{frame}

% ---------------------------------------------------------------
\begin{frame}{Resumen de la sesión 19}
  \begin{resumenbox}
    \begin{itemize}
      \item La revisión es una \textbf{parte fundamental} del proceso
            matemático, no un extra opcional.
      \item Hay 4 tipos de errores: \textbf{matemáticos},
            de datos, de representación y de comunicación.
      \item La \textbf{rúbrica} permite autoevaluar
            objetivamente con criterios claros.
      \item La \textbf{coevaluación} con el protocolo
            «Estrella + Deseo» da feedback constructivo.
      \item Cada corrección debe documentarse:
            original → problema → versión corregida.
    \end{itemize}
  \end{resumenbox}
  \vspace{6pt}
  \begin{center}
    {\small\color{AlmAzulM}\faArrowCircleRight\enskip
    \textbf{Sesión 20 (última):}
    Exposición final y cierre del proyecto}
  \end{center}
\end{frame}

\end{document}
